% Separate primary treatment: Rodkinson's printed-copyright-1900 translation.
% Not a reconstruction or discharge of Abraham--Gabbay--Schild HPL18 (2018).
\section{Talmudic inference: comparison, limitation and objection}
\label{sec:p18-talmudic-inference-principles}

An argument from a less demanding case to a more demanding one needs an
identified respect in which the second case is more demanding. It also
needs a reason for transferring the consequence, a limit on what may be
transferred, and an account of exceptions. The discussion of animal
damage in the Mishnah and Gemara supplies actual disputes about these
requirements, not merely a name for an argument form.

This entry follows Chapter II, printed pages 30--56, of Michael L.
Rodkinson's English translation of \emph{Baba Kama}, Volume II (X),
whose copyright page bears 1900
\parencite[pp.~30--56]{RodkinsonBabaKamaCopyright1900}.
Its central passage is Mishnah V and the following Gemara, pages 50--54.
That argument sequence is the focal reconstruction; identified neighboring
passages supply additional worked comparisons. This is not a line-by-line
reconstruction of every case in Chapter II.
The familiar terms \emph{qal va-homer} and \emph{dayo} designate,
respectively, the stronger-case inference and the limiting principle
rendered here as an a fortiori inference and a sufficient result.
The spelling \emph{Tarphon} is retained from this translation.
The later Abraham--Gabbay--Schild \emph{Principles of Talmudic Logic}
is a separate modern treatment, not the source of the reconstruction
given here \cite{AbrahamGabbaySchild2018PrinciplesMetadata}.

The propositions about liability, ritual categories and scriptural
authority below belong to the translated traditional arguments.
They are not current religious rulings, medical statements, or legal
advice. Describing how a passage reasons neither establishes its
normative premises independently nor recommends its rules for a
present dispute.

\subsection{Identify the textual layer before reconstructing the inference}
\label{sec:p18-talmudic-textual-layer}

The translation distinguishes a \emph{Mishna}, a \emph{Boraitha}, and
the surrounding Gemara's questions and replies. These are different
roles within its discussion: a stated teaching, a further cited
teaching, and reasoning about their interpretation and consequences.
The participants can agree on a cited rule while disagreeing about
which inference it licenses or which case it describes.

Rodkinson also declares editorial conventions in his explanatory
remarks. Questions are followed immediately by their answers.
Parentheses can carry Rashi's explanation, and square brackets carry
other commentary under the convention he states. Where the original
offers certain alternative formulations, he translates only the second.
He does not supply the original page numbers in this English edition
\parencite[Explanatory Remarks, p.~ii]{RodkinsonBabaKamaCopyright1900}.
Accordingly, an explanatory parenthesis is not silently promoted to
the unchanged words of the Mishnah, and an English passage containing
one alternative is not evidence that the underlying tradition had
only that alternative.

The reconstruction preserves the translated arguments and their
objections in ordinary prose. The comparison table below is an
editorial aid to keeping the premises straight; it is not a calculus,
an original table in the Talmud, or a complete modern axiomatization.

\subsection{The actual comparison: cause, place and amount}
\label{sec:p18-talmudic-damage-premises}

The source's categories do not mean simply that one animal is good and
another bad. The \emph{foot} concerns characteristic damage in walking;
the \emph{tooth} concerns characteristic consumption; the \emph{horn}
includes the specified goring, pushing, biting, lying down and kicking
cases. The chapter distinguishes ordinary from unusual behavior,
direct bodily action from transmitted force, and an animal's changed
status after warning. Those distinctions determine which comparison
is proposed.

Mishnah V takes a horn-type damage case with the half-damage rule on
public ground. For damage on the injured party's premises, Tarphon
argues for full damages, while the rabbis retain half. The shared
comparison with tooth and foot is:
\begin{center}
\begin{tabular}{p{0.31\linewidth}p{0.22\linewidth}p{0.36\linewidth}}
\hline
Cause in the compared case & Public ground & Injured party's premises\\
\hline
Tooth or foot, ordinary damage & No damage liability & Full damages\\
Horn in the specified half-damage case & Half damages &
Half, or full under Tarphon's proposed inference\\
\hline
\end{tabular}
\end{center}
The last entry is the disputed conclusion, not an additional agreed
premise \parencite[pp.~50--51]{RodkinsonBabaKamaCopyright1900}.

This table concerns the damage liability being compared in the
argument. It does not erase other consequence types. Earlier in
the chapter, an animal's consumption on public ground can still
raise a separate question about payment for benefit, even where
the damage rule is inapplicable. Likewise, the distinction between
non-vicious and vicious horn cases is not removed by recording a
single half-damage comparison
\parencite[pp.~37--40, 47--50]{RodkinsonBabaKamaCopyright1900}.
In the table, absence of damage liability therefore must not be read as
absence of every possible payment on any ground.

\subsection{Two orientations of the a fortiori argument}
\label{sec:p18-talmudic-qal-vahomer}

Tarphon's first argument compares places while holding a cause
category available for calibration. Public ground receives the
more lenient treatment for tooth and foot, yet horn damage there
already incurs half. On the injured party's premises, where tooth
and foot receive the stricter full-damage treatment, he argues that
horn damage should receive full damages as well.

The stated facts alone do not contain the disputed result. The
inferential move is that the stricter treatment of the place in the
accepted tooth/foot cases should carry over to horn. Its warrant is
not ordinary subtraction: the change from no damage liability to
full damages is not treated as an instruction to add a whole damage
amount to every existing half-damage amount. Neither a multiplier
nor a generally applicable arithmetic formula is stated here.

The rabbis reply with the limiting principle: the result established
by this inference is to match the originating case. Taking the
horn case on public ground as the originating liability comparison,
its half damages limit the derived horn result on private premises.
Their objection accepts the relevance of inference while denying
the increase Tarphon seeks. It is not the claim that no comparison
can ever be used.

Tarphon then reformulates. He says that he is not deriving horn from
horn, but horn from foot: where tooth and foot have the lenient
treatment on public ground, horn has the stricter treatment; therefore
horn should be at least as strictly treated where tooth and foot
already have full-damage liability. This second orientation compares
cause categories, using the public-ground contrast to establish their
ordering and carrying it to the injured party's premises.

The rabbis repeat the same limiting response. Re-describing the
comparison does not, for them, remove the limitation on the conclusion
\parencite[pp.~50--51]{RodkinsonBabaKamaCopyright1900}.
The disagreement thus includes the identification of the originating
case and which features of a multi-premise comparison the limit
controls. An analyst who writes only a stronger case and a stronger
result loses precisely the issue being disputed.

\subsection{Dayo and the dispute about an inference's remaining work}
\label{sec:p18-talmudic-dayo}

The Gemara next asks whether Tarphon ignores a limiting principle
which it treats as scripturally grounded. Its cited comparison is
the Miriam passage: the case of a father's rebuke has a seven-day
consequence; an argument from that case to the divine rebuke proposes
fourteen days; the limiting principle retains seven. These numbers
are part of this cited argument. The comparison does not prove that
every stronger case has exactly twice the consequence of a weaker one.

The reply distinguishes reducing an inference's proposed increase
from removing all of its independent contribution. In the account
ascribed to Tarphon, the divine-rebuke consequence was not separately
given as the starting rule. The comparison supplies a consequence,
and the limitation reduces it to seven; useful inferential work
remains. In the horn case, by contrast, half damages are already
given for the target setting in his account. The argument's new
contribution would be the additional half. Limiting its result back
to half removes all of that contribution. Tarphon is described as
rejecting the limitation in that latter circumstance, not as
rejecting it universally.

The translation's sentence introducing this exception is awkwardly
worded. The ensuing horn and Miriam illustrations identify the issue
as whether applying the limiting principle leaves any added result;
they do not establish a rule that an a fortiori inference refutes
itself. That contextual reading governs the reconstruction here.

The rabbis dispute the textual premise of the contrast. They point
to the stated seven-day exclusion in the Miriam passage.
Tarphon's response treats that statement itself as indicating the
limiting principle. The discussion then assigns work to the further
report of Miriam's actual exclusion: one reading uses it to support
the limitation, while the other uses it to show that the principle
extends beyond a special concern for Moses' honor
\parencite[pp.~51--52]{RodkinsonBabaKamaCopyright1900}.

Three questions must therefore remain separate. What result does
the stronger-case comparison initially propose? What originating
case limits that result? What has already been independently
established for the target? An inference can add a conclusion while
adding no increase in amount; conversely, an increase can be the
only new consequence at issue. The dispute about that distinction
depends on interpreting the cited text, not merely on choosing
between seven and fourteen.

An earlier use of the sufficient-result principle concerns the
source of payment, not only its amount. In the transmitted-force
discussion, a proposed liability is linked to horn, and that
originating category is invoked to constrain payment from the
tort-feasor's body rather than from another estate. The objection
that Tarphon does not accept the limitation is answered by
distinguishing the circumstances in which his rejection was
asserted \parencite[p.~35]{RodkinsonBabaKamaCopyright1900}.
The consequence being limited can thus contain more than a
single number. One must identify the payment mode as well as
the amount before claiming that two results are equivalent.

\begin{hbwarning}{def:p18-talmudic-dayo-strength}{A restriction and its exception are not a universal recipe}
The account of Tarphon's exception is one position in the translated
dispute. It is not a generally accepted permission to disregard a
limiting rule whenever an argument would otherwise become redundant.
Nor is the rabbis' reply represented as a theorem of numerical
monotonicity. Each position has its own claims about textual warrant,
the originating comparison and what the argument contributes.
\end{hbwarning}

\subsection{Proposed reverse comparisons and textual blockers}
\label{sec:p18-talmudic-blocked-comparisons}

The following discussion does not let the first comparison establish
a universal ordering of every cause category in every respect.
It tests other proposed consequences against explicit wording. The
sequence matters: each proposal looks attractive under one comparison,
but a cited restriction blocks its extension
\parencite[pp.~52--53]{RodkinsonBabaKamaCopyright1900}.

First, tooth and foot are proposed to incur liability on public
ground. The proposed comparison starts with horn: horn incurs only
half on private premises and still incurs half on public ground.
If tooth and foot incur full on private premises, why should they
be entirely exempt on public ground? The response reads the
other person's field in the cited passage as private, not public
ground. When a further objection proposes only half damages rather
than full, the reply assigns the half-damage wording to horn,
not to every cause of damage.

Second, tooth and foot are proposed to pay only half on private
premises. Horn is liable even on public ground, but pays only half
on private premises; should tooth and foot, exempt on public ground,
not receive the more lenient private-ground treatment? The reply
uses the restitution language to require the complete payment in
the tooth/foot case. The same comparison table has supported a
proposal of increasing one consequence and a proposal of decreasing
another, depending on which feature is treated as the ordering
dimension. The text-specific replies prevent these proposals from
being interchangeable applications of one numerical rule.

Third, horn is proposed to be exempt on public ground. Tooth and
foot incur full on private premises but are exempt in public;
horn incurs only half in private, so an inference toward the
more lenient public result is proposed. The additional division
wording is interpreted as preserving horn's liability in the
public-ground case.

The discussion then compares different consequence types in human
and animal injury cases. It proposes atonement payment for the
human case because a human is liable for the additional injury
payments while an ox, not liable for those payments, can occasion
atonement payment. The reply confines the cited atonement wording
to the specified ox case. A converse proposal would transfer the
additional injury payments to an ox, using its atonement consequence
as the stronger feature. The response cites the wording about a
person injuring a neighbor as excluding the ox from that transfer.
These are accounts of the source's proposed and rejected inferences,
not directions about present injury compensation.

Two lessons follow at the interpretive level. Being stricter in one
consequence does not establish a global order of subjects for all
consequences. Also, a blocked inference can be blocked by the scope
of an authoritative rule without demonstrating that the comparison
is meaningless in every setting. A reconstruction must record
the particular transfer proposed and the wording used against it.

\subsection{More frequent is not automatically stronger in the relevant respect}
\label{sec:p18-talmudic-frequency-countercase}

An earlier dispute asks when repeated goring establishes the
animal's changed status. The positions differ over three warnings
or incidents and over their distribution across days.
Meir's proposed stronger-case argument is that if three gorings
separated by long intervals suffice, three within one day should
suffice all the more. The feature taken as stronger is the
concentration of the incidents in time
\parencite[pp.~47--48]{RodkinsonBabaKamaCopyright1900}.

The reply gives a traditional ritual countercomparison. In the
cited category involving a woman's discharge, events on successive
days and repeated events within one day do not have the same
classification or consequence. More events within a shorter interval
therefore do not supply the proposed more-severe result in that
comparison. The inference from higher frequency to the stronger
consequence has met an example in which the number of days,
not merely the number of events, controls the classification.
Nothing here is a current medical account of the condition.

Meir answers by treating the ritual case as an explicitly marked
exception. The quoted wording is used to confine that case rather
than let it govern the animal comparison; the discussion also
contrasts the male and female ritual classifications.
The objection and response must both be preserved. The countercase
challenges a general warrant based on frequency. The exception
claim, if justified in the argument's traditional framework,
limits the countercase's reach. Merely calling it irrelevant
without its textual reason would omit the response.

The same section asks whether the days pertain to the ox's
classification alone or also to notification of the owner. Three
sets of witnesses arriving in a single day can attest three
incidents on different days. Those are three distinct temporal
objects: incidents, testimony and notification. The passage tests
the interpretations against the case of collusive witnesses, and
replies distinguish knowledge of the other witnesses' intended
testimony, signals between sets, joint arrival, and knowledge of
the owner rather than identification of the particular ox
\parencite[pp.~48--50]{RodkinsonBabaKamaCopyright1900}.
Counting three appearances does not settle which of those objects
the condition quantifies over.

\subsection{Similarity can conceal a difference that the argument needs}
\label{sec:p18-talmudic-benefit-loss}

The chapter's discussion of benefit supplies a concrete challenge
to analogical support. It asks whether an occupant who uses a
neighbor's court without the owner's knowledge must pay rent.
Before adducing a precedent, it separates the occupant's benefit
from the owner's loss. The difficult case is one in which the
occupant needs the space, while the owner was not offering it
for rent.

A proposed answer appeals to the rule requiring payment for the
benefit obtained when an animal consumes another person's food.
Rabha objects that this teaching concerns benefit to one party
with damage to another, whereas the occupancy question was posed
as benefit without loss. The analogy's shared feature, benefit,
does not by itself remove that difference. A bracketed reply in
the translation argues that food placed on public ground can be
treated as relinquished, seeking to restore the no-loss comparison;
its status as commentary is retained rather than attributed
unqualifiedly to the original speaker
\parencite[pp.~40--41]{RodkinsonBabaKamaCopyright1900}.

The later replies do not simply repeat the two labels. A cited
rent case is distinguished because a new building's walls can
be soiled by use. An apparent inconsistency between rent-free
occupation and payment to a subsequent owner is distinguished
by whether the occupant had intended to pay. Further explanations
point to the benefit of preventing a house's deterioration and
to the maintenance performed by an occupant. A storage case is
used to ask when the two reasons would yield different assessments
\parencite[pp.~41--42]{RodkinsonBabaKamaCopyright1900}.

Thus the operative premises concern which benefit, which loss,
whose knowledge, and what use of the property is at issue.
Calling both cases unauthorized use does not resolve those
questions. This is a source-native example of reconstructing
an analogy's material respects and testing its fit, not an
assertion that benefit and loss have been assigned a modern
utility function.

\subsection{Interpretive alternatives and their discriminating cases}
\label{sec:p18-talmudic-causal-classification}

The treatment of transmitted force gives another kind of
reasoning work. An animal steps on one vessel, whose fragments
then break a second; cocks' wings can break an object directly
or produce wind that does so. The discussion distinguishes
direct action from the gravel-type case and records Summachus'
different full-damage position. It then asks whether a vessel
which starts rolling under an animal's action and breaks later
is classified by the originating action or by the point of
breakage
\parencite[pp.~30--33]{RodkinsonBabaKamaCopyright1900}.

A cited hopping or water-pail case is offered as evidence for
an answer. The reply tests whether that teaching instead belongs
to a different position about gravel, whether primary and
secondary force differ, and whether the cock's action remained
continuous until the breakage. Each alternative changes what
the purported supporting case would show. A case that is
compatible with an alternative explanation does not uniquely
settle the proposed classification.

Similarly, the dog which carries a cake and a burning coal
to a barn is examined under competing accounts of fire as
an arrow and fire as property. The replies distinguish placing
the coal from flinging it, the place first struck from the
later spread, and a barrier that initially prevented spread
from an opportunity to repair it after it fell. They also
separate liability for the cake, the first place reached by
the coal and the barn
\parencite[pp.~43--46]{RodkinsonBabaKamaCopyright1900}.
Those distinctions are the argument's supplied causal and
classificatory premises. The metaphor of an arrow is not,
without them, a complete causal model.

The source's literal attributions and alternative readings
remain important here. In the storage example, the translation
reports Rashi's explanation and then expressly offers the
translator's different reading, leaving the choice open
\parencite[pp.~41--42]{RodkinsonBabaKamaCopyright1900}.
That alternative is not treated as an adjudicated correction
of Rashi or as proof that one entire causal doctrine is accepted.

\subsection{An unresolved question is not a negative conclusion}
\label{sec:p18-talmudic-open-questions}

Several questions are explicitly left unanswered or undecided.
The chapter asks whether gravel produced in an unusual manner
changes the fraction of damages, whether repeated gravel damage
changes status, and how to classify a kick where gravel cannot
be avoided. It also considers an extraordinary tail movement
\parencite[pp.~35--36]{RodkinsonBabaKamaCopyright1900}.
The alternatives and the unresolved standing belong to the
argument; replacing them by a universal rule would add a result
the translated discussion has not supplied.

The provoked-dog case makes the distinction particularly clear.
A proposed answer uses the description of a non-vicious ox as
one which does not gore when children pat it. From the failure
of that description, the proposal seeks liability for the
provoked event. Abayi challenges the reading: the teaching
could instead mean that the ox has lost the fully non-vicious
classification, without asserting liability for that particular
goring. The question is left undecided
\parencite[p.~50]{RodkinsonBabaKamaCopyright1900}.
Failure of a criterion for one classification does not itself
establish a different event-specific consequence.

This is not a rejection of inference from the stated rule.
It is a challenge to what the rule states and to the unstated
link between classification and the consequence sought.
Suspending that proposed answer therefore neither establishes
liability nor establishes exemption. The uncertainty must
remain attached to the question actually left open.

\subsection{Preserve the consequence type in a transferred comparison}
\label{sec:p18-talmudic-consequence-types}

The discussion also tests whether a foot-caused death on the
injured party's premises can support an atonement-payment consequence.
The proposed comparison with horn uses repeated habitual goring;
an objection distinguishes intention to harm in horn damage from
the foot's ordinary movement. The passage then examines a teaching
attributed to Tarphon, which imposes the full atonement consequence
for the specified unauthorized-entry and fatal-goring case.
It explains the inference using a half-atonement premise attributed
to Jose the Galilean and a foot-case comparison
\parencite[pp.~53--54]{RodkinsonBabaKamaCopyright1900}.

The final objection asks whether the proposed origin could instead
be the foot's ordinary property damage. That would not settle the
death case: the foot's ordinary tendency to damage objects does not
by itself establish the same tendency in killing. The reply therefore
locates the comparison in the atonement case itself, not merely in
ordinary damage. The whole line of reasoning is conditional on the
particular teachings invoked in the passage; it is not an independent
proof of their authority.

This protects two distinctions which the initial payment table
cannot express by itself. Damage payment and atonement payment
are different conclusions. Also, a classification describing
characteristic movement need not describe a characteristic
fatal result. Identical words such as full or half do not make
the two consequence types interchangeable.

\subsection{Comparison with a modern argument reconstruction}
\label{sec:p18-talmudic-modern-boundary}

The premise and warrant analysis of
\cref{sec:p18-argument-schemes-critical-questions} can help identify
what is stated, what is attributed implicitly, and what remains
unestablished. In the present cases, the ordering premise concerns
a specific cause, place, frequency or consequence. The transfer
warrant and the response invoking an exception are further
commitments. This is an editorial comparison between reasoning tasks,
not a claim that the Talmud states that modern scheme framework.

Likewise, the strict and defeasible rules of
\cref{sec:p18-is1996} are supplied inputs to a modern argument system.
Encoding a participant's proposed comparison in that system would
require deciding its premises, its strength, what counts as a
counterargument, and what priority an explicit textual restriction
receives. The availability of the framework does not answer those
interpretive questions. No preference ordering or inference calculus
for the whole traditional discussion is inferred from the case labels.

The legal-method distinction in \cref{def:p11-rat113-legal-logic}
also applies as a comparison of analytical questions. Whether a
conclusion follows from a stipulated rule is different from which
source is authoritative, how its words are interpreted, and which
case facts satisfy it. Here the translated debate supplies examples
of all those dependencies; it does not supply modern jurisdictional
authority, a psychological experiment, or an executable decision
procedure.

\begin{hbmethod}{def:p18-talmudic-reading-procedure}{Reconstructing one of the translated arguments}
This is an editorial reading procedure grounded in the worked cases,
not a newly attributed traditional rule or a procedure for issuing
a religious ruling.
\begin{enumerate}
\item Identify the exact layer and speaker: the teaching, a quoted
      comparison, the surrounding objection, or a translator's explanation.
\item State the source case and target case, holding the relevant
      conditions explicit. Separate place, manner of damage, status
      after warning, benefit, loss and consequence type.
\item Identify the respect in which one case is claimed to be stronger.
      Do not replace a stated comparison by an invented global severity
      scale or an arithmetic multiplier.
\item State the proposed transfer and identify the originating case
      used by the limiting response. If the limitation is challenged,
      preserve the speaker's condition on that challenge.
\item Record the actual defeater or reply: a textual restriction,
      a material difference, an alternative reading of the cited
      teaching, or an exception claim with its stated warrant.
\item Distinguish a supported answer, a contested answer, and a question
      the passage leaves undecided. Keep the uncertainty on that
      question rather than conclude its contrary.
\end{enumerate}
\end{hbmethod}

\subsection{Subject and task lookup}
\label{sec:p18-talmudic-lookup}

Use this treatment for the inferential work of comparing a source
case with a target, identifying what a limiting principle limits,
testing an analogy against a material difference, and following
an objection and reply through their actual premises. The complete
worked Tarphon case requires both orientations of the comparison,
the rabbis' repeated limitation, and the Gemara's dispute about
what inferential work remains.

A successful reading preserves the source-specific reason for the
ordering and the precise conclusion type. A failed reading infers
full damages merely because the word stronger appears, treats a
ritual exception as an unrestricted counterexample without its
reply, or turns loss of a classification into event-specific
liability without the missing link.

The scope is these arguments in the identified English translation.
It does not present every traditional hermeneutic principle, every
original textual variant, or the full modern Abraham--Gabbay--Schild
presentation. No formal soundness, completeness, decidability or
proof of a religious premise is claimed.
